\documentclass[border=6pt]{standalone}
\usepackage[T1]{fontenc}
\usepackage[utf8]{inputenc}
\usepackage[brazilian]{babel}
\usepackage[helvratio=0.9]{newpxtext}
\usepackage{newpxmath}
\usepackage[scaled=0.94]{inconsolata}
\usepackage{fontawesome5}
\usepackage{tikz}
\usepackage{pgfplots}
\pgfplotsset{compat=1.17}
\usetikzlibrary{arrows.meta,positioning,calc,fit,backgrounds,shapes.geometric,shapes.misc,decorations.pathreplacing}
\definecolor{corprimaria}{HTML}{2C7A73}
\definecolor{corprimariaescura}{HTML}{1F5C56}
\definecolor{cortinta}{HTML}{17252D}
\definecolor{corcinza}{HTML}{5C6870}
\definecolor{corsuave}{HTML}{9AA5AC}
\definecolor{cortom}{HTML}{EDF4F3}
\definecolor{corpapel}{HTML}{F7F9F9}
\definecolor{sevcritica}{HTML}{8E1B1B}
\definecolor{sevalta}{HTML}{BF4E15}
\definecolor{sevmedia}{HTML}{A8790A}
\definecolor{sevbaixa}{HTML}{3A6B8A}
\definecolor{sevinformativa}{HTML}{6B7280}
\definecolor{sitatende}{HTML}{2E7D4F}
\tikzset{
  every picture/.style={font=\sffamily\footnotesize,text=cortinta},
  caixa/.style={draw=corsuave,line width=0.5pt,rounded corners=2pt,fill=white,inner sep=6pt,align=center},
  caixaforte/.style={caixa,draw=corprimaria,line width=0.9pt,fill=cortom},
  seta/.style={-{Stealth[length=5pt,width=4pt]},line width=0.6pt,draw=corcinza},
  setaforte/.style={seta,draw=corprimaria,line width=0.9pt},
  rotulo/.style={font=\sffamily\scriptsize,text=corcinza,align=center,inner sep=1.5pt},
  titulo/.style={font=\sffamily\bfseries\footnotesize,text=corprimaria},
}
\pgfplotsset{
  relatorio/.style={
    width=13.5cm,height=6.6cm,
    axis lines*=left,axis line style={draw=corsuave},
    tick style={draw=corsuave},tick label style={font=\sffamily\footnotesize,text=corcinza},
    label style={font=\sffamily\footnotesize,text=corcinza},
    ymajorgrids,grid style={draw=corsuave!35,line width=0.3pt},
    legend style={draw=none,fill=none,font=\sffamily\footnotesize,legend cell align=left},
    every node near coord/.append style={font=\sffamily\scriptsize\bfseries,text=cortinta},
    /pgf/number format/.cd,use comma,1000 sep={.},assume math mode=true,
  },
}

\newcommand{\nDimensoes}{?}
\newcommand{\nRefutacoes}{?}
\newcommand{\nRefutPorDim}{?}
\newcommand{\nLacunas}{?}
\newcommand{\nAchados}{?}
\newcommand{\nCriticos}{?}
\newcommand{\nAltos}{?}
\newcommand{\nMedios}{?}
\newcommand{\nBaixos}{?}
\newcommand{\nInformativos}{?}
\newcommand{\Prefixo}{LGPD}
\newcommand{\TituloLista}{As dimensões analisadas}
\newcommand{\ListaDimensoes}{%
  1.\ inventário de dados pessoais & 4.\ direitos, retenção e eliminação & 7.\ dados de terceiros e sigilo\\
  2.\ operadores e transferência internacional & 5.\ segurança e incidentes & 8.\ lojas e declarações públicas\\
  3.\ transparência, bases legais e aceite & 6.\ inteligência artificial & 9.\ versão web}
\newcommand{\num}[1]{{\bfseries\LARGE\color{corprimaria}#1}}
\newcommand{\etapa}[6]{%
  \node[#1,minimum width=2.62cm,minimum height=4.15cm,anchor=north,inner sep=0pt] (#2) at (#3,0) {};
  \node[anchor=north,text width=2.36cm,align=center,inner sep=0pt] (#2t) at ([yshift=-0.42cm]#2.north) {\textbf{#4}};
  \node[anchor=north,inner sep=0pt] (#2n) at ([yshift=-1.22cm]#2.north) {\num{#5}};
  \node[anchor=north,text width=2.36cm,align=center,inner sep=0pt,font=\sffamily\scriptsize] at ([yshift=-1.95cm]#2.north) {#6};
}
\begin{document}
\hyphenpenalty=10000\exhyphenpenalty=10000
\begin{tikzpicture}
\tikzset{
  passo/.style={font=\sffamily\bfseries\scriptsize,text=white,fill=corprimaria,circle,inner sep=1.6pt,minimum size=13pt},
}
\etapa{caixa}{e1}{0}{Análise\\por dimensão}{\nDimensoes}{dimensões em paralelo, um dossiê por dimensão}
\etapa{caixa}{e2}{3.02}{Verificação\\independente}{\nDimensoes}{verificações, uma por dossiê, com correção no próprio arquivo}
\etapa{caixa}{e3}{6.04}{Refutação\\adversarial}{\nRefutacoes}{refutações de achados graves, até \nRefutPorDim\ por dimensão}
\etapa{caixa}{e4}{9.06}{Crítico de\\completude}{\nLacunas}{lacunas, cada uma em uma apuração complementar}
\etapa{caixaforte}{e5}{12.08}{Registro\\mestre}{\nAchados}{achados únicos, de \Prefixo-01 a \Prefixo-\nAchados}
\foreach \i in {1,...,5} {\node[passo] at (e\i.north) {\i};}
\foreach \a/\b in {e1/e2,e2/e3,e3/e4,e4/e5} {\draw[setaforte] (\a.east) -- (\b.west);}
\draw[draw=corprimaria!45,line width=0.4pt] ([xshift=6pt,yshift=1.16cm]e5.south west) -- ([xshift=-6pt,yshift=1.16cm]e5.south east);
\node[anchor=south,font=\sffamily\scriptsize,align=center,inner sep=0pt] at ([yshift=7pt]e5.south) {%
  {\color{corcinza}\textbf{\nCriticos}} críticos\enspace{\color{sevalta}\textbf{\nAltos}} altos\\[1pt]
  {\color{sevmedia}\textbf{\nMedios}} médios\enspace{\color{sevbaixa}\textbf{\nBaixos}} baixos\\[1pt]
  {\color{sevinformativa}\textbf{\nInformativos}} informativos};
\draw[decorate,decoration={brace,mirror,amplitude=4pt},draw=corsuave,line width=0.5pt]
  ([yshift=-5pt]e1.south west) -- ([yshift=-5pt]e1.south east);
\node[rotulo,anchor=north west,align=left] at ([xshift=-2pt,yshift=-14pt]e1.south west) {%
  \textbf{\color{cortinta}\TituloLista}\\[2pt]
  \begin{tabular}{@{}l@{\hspace{1.6em}}l@{\hspace{1.6em}}l@{}}
  \ListaDimensoes
  \end{tabular}};
\end{tikzpicture}
\end{document}
